\chapter{Closing}

La fase di chiusura si è aperta con la richiesta a Giovanni Marchetti dell'accettazione formale della soluzione. L'11 Maggio 2026 Marco Venturi ha presentato allo sponsor tutti i deliverable del MVP e Giovanni, dopo aver testato personalmente una partita completa, ha firmato il documento di accettazione senza richiedere modifiche. L'installazione sui server di produzione non è stata a carico di PlayHeritage Labs: il deploy finale è stato gestito dal team tecnico della community, utilizzando le Docker image e le istruzioni di deployment preparate da Andrea Conti.

La documentazione prodotta durante il progetto è raccolta integralmente nel Project Notebook --- il workspace Notion del progetto, alimentato fin dal primo giorno ---, di cui il Final Project Report, consegnato a Giovanni il 12 Maggio 2026, costituisce il documento di sintesi. Il report e il dettaglio delle attività di chiusura --- audit, fattori di successo, lezioni apprese e prospettive future --- sono riportati nel \allegatolink{Documentazione/Closing/Capitolo6-Closing.pdf}{documento di dettaglio Closing} allegato alla relazione.

\section{Audit Post-Implementazione}

Da un audit post-implementazione condotto insieme a Giovanni Marchetti è emerso che tutti gli obiettivi del Project Overview Statement sono stati raggiunti e i criteri di successo pienamente soddisfatti: progetto completato in 7 mesi nei tempi previsti (141 giorni lavorativi sul critical path); budget rispettato (€22.750 su €25.000, con un surplus di €2.250 destinato al supporto post-lancio); regole del Maraffone implementate e validate da Francesca Giuliani; latenza real-time media di 185ms, sotto il target di 500ms; e l'80\% degli utenti in grado di completare una partita senza assistenza.

\section{Lessons Learned e Chiusura}

I fattori che hanno reso possibile questo risultato sono stati il monitoraggio continuo e multilivello, la scelta di metodologie differenziate per sottosistema, la collaborazione costante dello sponsor e la competenza del team. Tra le lezioni apprese, le più rilevanti riguardano l'opportunità di prototipare le feature UI complesse già in fase di stima, di coinvolgere gli esperti di dominio non solo nello Scoping ma anche durante lo sviluppo, e di anticipare il testing cross-browser rispetto agli sprint finali. Il dettaglio completo dei fattori di successo e delle lezioni apprese è riportato nel \allegatolink{Documentazione/Closing/Capitolo6-Closing.pdf}{documento di dettaglio Closing}.

La soddisfazione del committente si è tradotta nell'interesse per sviluppi futuri della piattaforma --- app mobile nativa, sistema di tornei, modalità single-player contro AI --- da valutare nei mesi successivi al lancio. Il progetto si chiude raggiungendo gli obiettivi fissati e confermando che è possibile digitalizzare un gioco tradizionale preservandone le regole e lo spirito, applicando le metodologie di Project Management con flessibilità e pragmatismo.
